\chapter{Transpilação: Geração de Código C}\label{chap:transpilacao}

\section{Introdução}\label{sec:transp-intro}

Ao longo dos capítulos anteriores construímos o \emph{pipeline} clássico de
compilação: análise léxica e sintática, verificação de tipos, geração de
representação intermediária e, no capítulo seguinte, tradução para código de
máquina.  Existe, contudo, uma rota alternativa — amplamente usada na indústria
— que substitui a geração de código de máquina por uma etapa distinta:
traduzir o programa para outra \emph{linguagem de alto nível}.  Essa técnica
chama-se \textbf{transpilação} (ou \emph{source-to-source compilation}).

\subsection{Por que transpilar?}

Um transpilador não produz código de máquina: ele delega essa tarefa ao
compilador da linguagem-alvo.  Em troca, obtém-se:

\begin{itemize}
  \item \textbf{Portabilidade imediata.}  Basta que o compilador-alvo exista
    na plataforma; não é necessário implementar um \emph{backend} por
    arquitetura.
  \item \textbf{Otimizações do compilador-alvo.}  O compilador-alvo
    (e.g.~\texttt{gcc -O2}) aplica análise de fluxo, alocação de
    registradores e inlining de forma automática.
  \item \textbf{Simplicidade de implementação.}  Sem IRT, sem alocação de
    registradores, sem linearização de controle de fluxo.
  \item \textbf{Código gerado legível.}  O arquivo \texttt{.c} produzido pode
    ser inspecionado e depurado diretamente.
\end{itemize}

A contrapartida é que o transpilador não controla a qualidade do código final
e está sujeito às restrições de expressividade da linguagem-alvo.

\begin{center}
  \begin{tikzpicture}[
      node distance=1.0cm and 1.8cm,
      box/.style={draw, rounded corners, minimum width=2.8cm, minimum
        height=0.9cm, font=\small\sffamily, align=center},
      arr/.style={-Stealth, thick}
    ]
    \node[box] (ast)  {AST tipada};
    \node[box, right=of ast]  (c)     {Código C};
    \node[box, right=of c]    (bin)   {Binário};
    \draw[arr] (ast) -- (c)   node[midway, above, font=\tiny\sffamily]{transpilador};
    \draw[arr] (c)   -- (bin) node[midway, above, font=\tiny\sffamily]{\texttt{gcc}};
  \end{tikzpicture}
\end{center}

\subsection{TypeScript: um transpilador amplamente utilizado}

O exemplo mais proeminente de transpilador na indústria atual é o
\textbf{TypeScript}.  TypeScript é uma linguagem de programação desenvolvida
pela Microsoft que adiciona um \textbf{sistema de tipos estático} ao
JavaScript.  O compilador \texttt{tsc} recebe arquivos \texttt{.ts} e produz
arquivos \texttt{.js} válidos, executáveis em qualquer ambiente JavaScript
— navegadores, Node.js, Deno, etc.

A relevância de TypeScript ilustra vários aspectos importantes da
transpilação:

\begin{description}
  \item[Tipos apagados em tempo de execução.]  O código C, Java ou Rust
    gerado por um compilador carrega informação de tipo em tempo de execução
    (via RTTI ou layouts de memória tipados).  Em TypeScript, os tipos são
    \emph{completamente apagados} após a verificação: o \texttt{.js} gerado
    não contém nenhum rastro das anotações de tipo.  A segurança de tipos
    existe apenas em tempo de compilação.

  \item[Compatibilidade com a linguagem-alvo.]  TypeScript precisa gerar
    JavaScript válido para todas as versões-alvo (\texttt{ES5}, \texttt{ES2020},
    etc.).  Construções modernas como classes, destructuring e async/await são
    transpiladas para equivalentes mais antigos quando necessário.  Isso impõe
    restrições: TypeScript não pode gerar código JavaScript que expresse
    conceitos sem equivalente na linguagem-alvo.

  \item[Escala industrial.]  O compilador \texttt{tsc} é usado por dezenas
    de milhões de desenvolvedores.  Sua adoção exponencial (de 2012 a hoje)
    demonstra que a transpilação não é apenas um recurso pedagógico: é uma
    estratégia de engenharia viável para linguagens de produção.

  \item[Ecosistema de ferramentas.]  Transpiladores podem gerar
    \emph{source maps} — arquivos que mapeiam posições no código gerado de
    volta ao código-fonte original — permitindo que depuradores exibam o
    código TypeScript enquanto executam JavaScript.  Essa técnica é aplicável
    a qualquer transpilador.
\end{description}

Outros transpiladores notáveis incluem:
\begin{itemize}
  \item \textbf{CoffeeScript $\to$ JavaScript} (2009): precursor de TypeScript,
    oferecia sintaxe mais concisa.
  \item \textbf{Kotlin/JS $\to$ JavaScript}: o compilador Kotlin pode
    transpilar para JavaScript além de JVM bytecode.
  \item \textbf{Elm $\to$ JavaScript}: linguagem funcional com tipos algébricos
    transpilada para JavaScript.
  \item \textbf{Babel}: transpila JavaScript moderno (ES2022+) para JavaScript
    compatível com navegadores mais antigos.
  \item \textbf{\texttt{cfront}}: o compilador original de C++ (anos 1980)
    gerava código C; C++ era inicialmente uma extensão de C.
  \item \textbf{GHC}: o compilador de Haskell suporta um \emph{backend} C
    que gera código C portável como passo intermediário.
\end{itemize}

\subsection{Escopo deste capítulo}

Apresentamos geradores de código C para as três linguagens da disciplina —
TLine, TWhile e TImp —, em ordem crescente de complexidade.  Cada gerador
transforma diretamente a AST tipada em um arquivo \texttt{.c} válido que,
compilado com \texttt{gcc}, produz um executável com a mesma semântica do
interpretador apresentado nos capítulos anteriores.

% ============================================================================
\section{C como Linguagem-Alvo}\label{sec:c-target}
% ============================================================================

C é uma escolha natural para a linguagem-alvo neste contexto.  Seu modelo de
execução — memória plana, controle de fluxo estruturado, funções com escopo
local — mapeia-se diretamente às construções de TLine, TWhile e TImp.  A
biblioteca padrão (\texttt{stdio.h}, \texttt{stdlib.h}, \texttt{string.h})
provê as funções de I/O e alocação de memória necessárias sem dependências
externas.

\subsection{Representação de tipos}

A tabela~\ref{tab:tipo-c} resume como cada tipo das linguagens-fonte é
representado em C.

\begin{table}[htbp]
\centering
\begin{tabular}{lll}
\hline
\textbf{Tipo fonte} & \textbf{C (expressão / parâmetro)} & \textbf{C (variável local)} \\
\hline
\texttt{int}        & \texttt{long}    & \texttt{long v}         \\
\texttt{bool}       & \texttt{int}     & \texttt{int v}          \\
\texttt{string}     & \texttt{char *}  & \texttt{char v[4096]}   \\
\texttt{R} (record) & \texttt{R *}     & \texttt{R *v}           \\
\hline
\end{tabular}
\caption{Mapeamento de tipos para C.}\label{tab:tipo-c}
\end{table}

A distinção entre ``expressão'' e ``variável local'' para strings existe porque
parâmetros de função são \texttt{char *} (ponteiro para buffer existente),
enquanto variáveis locais precisam de um buffer próprio (\texttt{char[4096]})
para suportar leitura via \texttt{fgets}.  Records são sempre representados
como ponteiros para structs alocados no heap com \texttt{malloc}.

\subsection{A mônada de geração \texttt{CgM}}

Todos os três transpiladores compartilham a mesma estratégia de
implementação: uma \textbf{mônada de estado} que acumula linhas de código C
e mantém o ambiente necessário para as decisões de geração.

\begin{minted}{haskell}
data CgSt = CgSt
  { cgLines  :: [String]   -- linhas de código acumuladas
  , cgEnv    :: Map Var Ty -- tipos das variáveis em escopo
  , cgIndent :: Int        -- nível de indentação corrente
  }

type CgM a = StateT CgSt (ExceptT String Identity) a
\end{minted}

O tipo \texttt{CgM a} combina:
\begin{itemize}
  \item \texttt{StateT CgSt}: leitura e atualização do estado de geração;
  \item \texttt{ExceptT String}: possibilidade de falha com mensagem de erro;
  \item \texttt{Identity}: sem outros efeitos colaterais.
\end{itemize}

Duas operações primitivas sustentam toda a geração de código:

\begin{minted}{haskell}
-- Emite uma linha com o recuo corrente.
emit :: String -> CgM ()
emit line = modify $ \s ->
  s { cgLines = cgLines s ++ [replicate (cgIndent s * 4) ' ' ++ line] }

-- Executa m com um nível extra de indentação.
indented :: CgM () -> CgM ()
indented m = do
  modify $ \s -> s { cgIndent = cgIndent s + 1 }
  m
  modify $ \s -> s { cgIndent = cgIndent s - 1 }
\end{minted}

O ponto de entrada de cada transpilador executa a mônada e coleta as linhas:

\begin{minted}{haskell}
runCg :: CgM () -> Either String String
runCg m =
  case runIdentity (runExceptT (execStateT m initSt)) of
    Left err -> Left err
    Right st -> Right (unlines (cgLines st))
  where
    initSt = CgSt [] Map.empty 0
\end{minted}

% ============================================================================
\section{Geração de Código para TLine}\label{sec:codegen-tline}
% ============================================================================

TLine é a linguagem mais simples da família: declarações de variáveis,
atribuições e operações de I/O, sem controle de fluxo.  O programa
\texttt{TLine} é uma lista de comandos executados sequencialmente em um único
escopo.

\subsection{Estrutura do código gerado}

Todo programa TLine é traduzido para um único arquivo C com a seguinte
estrutura:

\begin{minted}{c}
#include <stdio.h>
#include <stdlib.h>
#include <string.h>

int main(void) {
    /* tradução dos comandos */
    return 0;
}
\end{minted}

O ponto de entrada do transpilador reflete exatamente essa estrutura:

\begin{minted}{haskell}
compileTLine :: TLine -> Either String String
compileTLine (TLine stmts) = runCg $ do
  emit "#include <stdio.h>"
  emit "#include <stdlib.h>"
  emit "#include <string.h>"
  blank
  emit "int main(void) {"
  indented $ do
    mapM_ genStmt stmts
    emit "return 0;"
  emit "}"
\end{minted}

\subsection{Geração de expressões}

A função \texttt{genExp} traduz uma expressão da AST para uma string C.
Ela recebe explicitamente o ambiente de tipos \texttt{env :: Map Var Ty},
necessário para inferir o tipo de variáveis em comandos de I/O:

\begin{minted}{haskell}
genExp :: Map Var Ty -> Exp -> String
genExp _   (EInt n)      = show n ++ "L"
genExp _   (EBool True)  = "1"
genExp _   (EBool False) = "0"
genExp _   (EString s)   = "\"" ++ escapeC s ++ "\""
genExp _   (EVar v)      = v
genExp env (e1 :+: e2)   = parens (genExp env e1 ++ " + "  ++ genExp env e2)
-- demais operadores binários são análogos
genExp env (Not e)        = "(!" ++ genExp env e ++ ")"
\end{minted}

Algumas decisões de representação merecem destaque:

\begin{itemize}
  \item Literais inteiros recebem o sufixo \texttt{L} para forçar o tipo
    \texttt{long} em C, prevenindo truncamentos em sistemas 32 bits.
  \item Literais booleanos são mapeados para \texttt{1}/\texttt{0}, pois C
    não tem tipo booleano nativo.
  \item Todos os operadores binários são envolvidos por parênteses (função
    \texttt{parens}), garantindo que a precedência da AST seja preservada
    independentemente da precedência de C.
  \item Strings literais passam pela função \texttt{escapeC} que escapa
    caracteres especiais (\texttt{"}, \texttt{\textbackslash}, \texttt{\textbackslash n}).
\end{itemize}

\paragraph{Regras de tradução de expressões.}

Definindo $\mathcal{E}_\Gamma\llbracket \cdot \rrbracket$ como a função de
tradução de expressões no ambiente de tipos $\Gamma$:

\begin{Definition}[Tradução de Expressões TLine]\label{def:tline-exp}
\begin{align*}
\mathcal{E}_\Gamma\llbracket n \rrbracket &= \mathtt{nL} \\
\mathcal{E}_\Gamma\llbracket \mathbf{true} \rrbracket &= \mathtt{1} \qquad
\mathcal{E}_\Gamma\llbracket \mathbf{false} \rrbracket = \mathtt{0} \\
\mathcal{E}_\Gamma\llbracket s \rrbracket &= \mathtt{\textquotedbl}\mathit{escapeC}(s)\mathtt{\textquotedbl} \qquad
\mathcal{E}_\Gamma\llbracket x \rrbracket = x \\
\mathcal{E}_\Gamma\llbracket e_1 \oplus e_2 \rrbracket
  &= \mathtt{(}\mathcal{E}_\Gamma\llbracket e_1 \rrbracket\ \mathit{opC}(\oplus)\ \mathcal{E}_\Gamma\llbracket e_2 \rrbracket\mathtt{)} \\
\mathcal{E}_\Gamma\llbracket \mathbf{not}\ e \rrbracket
  &= \mathtt{(!}\mathcal{E}_\Gamma\llbracket e \rrbracket\mathtt{)}
\end{align*}
onde $\mathit{opC}$ mapeia: $+{\mapsto}\mathtt{+}$,
$-{\mapsto}\mathtt{-}$, $*{\mapsto}\mathtt{*}$,
$/{}\mapsto\mathtt{/}$, $<{\mapsto}\mathtt{<}$,
$={\mapsto}\mathtt{==}$.
\end{Definition}

\subsection{Geração de comandos}

\paragraph{Declaração (\texttt{SDecl v t e}).}
Para tipos não-string, a tradução é uma declaração com inicializador direto.
Para strings, é necessária uma declaração de buffer seguida de
\texttt{strcpy}, pois C não suporta atribuição de arrays com \texttt{=}:

\begin{minted}{haskell}
genStmt (SDecl v t e) = do
  env <- gets cgEnv
  let rhs = genExp env e
  case t of
    TString -> do
      emit (localDecl v TString ++ ";")
      emit ("strcpy(" ++ v ++ ", " ++ rhs ++ ");")
    _ ->
      emit (localDecl v t ++ " = " ++ rhs ++ ";")
  setVar v t
\end{minted}

\paragraph{Atribuição (\texttt{SAssign v e}).}
Análoga à declaração, mas sem o tipo.  O tipo corrente de $v$ é consultado
em \texttt{cgEnv}:

\begin{minted}{haskell}
genStmt (SAssign v e) = do
  env <- gets cgEnv
  ty  <- getVarTy v
  let rhs = genExp env e
  case ty of
    TString -> emit ("strcpy(" ++ v ++ ", " ++ rhs ++ ");")
    _       -> emit (v ++ " = " ++ rhs ++ ";")
\end{minted}

\paragraph{Impressão (\texttt{SPrint e}).}
O formato de \texttt{printf} é escolhido pelo tipo inferido de $e$.  A função
auxiliar \texttt{inferTy :: Map Var Ty -> Exp -> Ty} determina o tipo de uma
expressão consultando \texttt{cgEnv} para variáveis:

\begin{minted}{haskell}
genStmt (SPrint e) = do
  env <- gets cgEnv
  let ce = genExp env e
  case inferTy env e of
    TInt    -> emit ("printf(\"%ld\\n\", (long)(" ++ ce ++ "));")
    TBool   -> emit ("printf(\"%s\\n\", (" ++ ce ++ ") ? \"true\" : \"false\");")
    TString -> emit ("printf(\"%s\\n\", " ++ ce ++ ");")
\end{minted}

\paragraph{Leitura (\texttt{SRead prompt v}).}
O prompt é impresso com \texttt{printf} e o valor é lido com \texttt{scanf}
(inteiros e booleanos) ou \texttt{fgets} (strings).  Para strings, o
caractere de nova linha inserido por \texttt{fgets} é removido:

\begin{minted}{haskell}
genStmt (SRead prompt v) = do
  env <- gets cgEnv
  emit ("printf(\"%s\", " ++ genExp env prompt ++ ");")
  ty <- getVarTy v
  case ty of
    TInt    -> emit ("scanf(\"%ld\", &" ++ v ++ ");")
    TBool   -> emit ("scanf(\"%d\",  &" ++ v ++ ");")
    TString -> do
      emit ("fgets(" ++ v ++ ", 4096, stdin);")
      emit (v ++ "[strcspn(" ++ v ++ ", \"\\n\")] = '\\0';")
\end{minted}

\paragraph{Regras de tradução de comandos.}

\begin{Definition}[Tradução de Comandos TLine]\label{def:tline-stmt}
\begin{align*}
\mathcal{S}_\Gamma\llbracket \mathbf{var}\ v : \tau = e \rrbracket
  &= \begin{cases}
       \mathtt{char\ }v\mathtt{[4096];\ strcpy(}v\mathtt{,\ }\mathcal{E}_\Gamma\llbracket e\rrbracket\mathtt{);}
         & \tau = \mathtt{string} \\
       \mathit{cTy}(\tau)\ v\ \mathtt{=}\ \mathcal{E}_\Gamma\llbracket e\rrbracket\mathtt{;}
         & \text{c.c.}
     \end{cases} \\[6pt]
\mathcal{S}_\Gamma\llbracket v \mathrel{:=} e \rrbracket
  &= \begin{cases}
       \mathtt{strcpy(}v\mathtt{,\ }\mathcal{E}_\Gamma\llbracket e\rrbracket\mathtt{);}
         & \Gamma(v) = \mathtt{string} \\
       v\ \mathtt{=}\ \mathcal{E}_\Gamma\llbracket e\rrbracket\mathtt{;}
         & \text{c.c.}
     \end{cases} \\[6pt]
\mathcal{S}_\Gamma\llbracket \mathbf{print}\ e \rrbracket
  &= \mathit{printf}_C\bigl(\mathit{ty}_\Gamma(e),\ \mathcal{E}_\Gamma\llbracket e\rrbracket\bigr)
\end{align*}
onde $\mathit{printf}_C(\mathtt{int},\ c) = \mathtt{printf("\%ld\textbackslash n",\ (long)(}c\mathtt{));}$,
$\mathit{printf}_C(\mathtt{bool},\ c) = \mathtt{printf("\%s\textbackslash n",\ (}c\mathtt{)\ ?\ "true"\ :\ "false");}$,
$\mathit{printf}_C(\mathtt{string},\ c) = \mathtt{printf("\%s\textbackslash n",\ }c\mathtt{);}$.
\end{Definition}

\subsection{Exemplo: soma de dois inteiros}

O programa TLine:
\begin{minted}{text}
var a : int = 0;
var b : int = 0;
read "a: " a;
read "b: " b;
print a + b;
\end{minted}

é transpilado para:

\begin{minted}{c}
#include <stdio.h>
#include <stdlib.h>
#include <string.h>

int main(void) {
    long a = 0L;
    long b = 0L;
    printf("%s", "a: ");
    scanf("%ld", &a);
    printf("%s", "b: ");
    scanf("%ld", &b);
    printf("%ld\n", (long)((a + b)));
    return 0;
}
\end{minted}

% ============================================================================
\section{Geração de Código para TWhile}\label{sec:codegen-twhile}
% ============================================================================

TWhile estende TLine com laços \texttt{while}, condicionais
\texttt{if}/\texttt{else} e operadores booleanos \texttt{\&\&} e \texttt{||}.
A tradução desses construtores para C é direta, mas exige atenção especial ao
\textbf{escopo léxico} de variáveis declaradas dentro de blocos.

\subsection{Escopo léxico e \texttt{withBlock}}

Em TWhile (e em C), uma variável declarada dentro de um bloco \texttt{while}
ou \texttt{if} não é visível fora dele.  O código C gerado já respeita esse
escopo por meio das chaves \texttt{\{ \}}.  Entretanto, o ambiente de tipos
\texttt{cgEnv} é um estado global da mônada; sem um mecanismo de restauração,
variáveis declaradas dentro de um bloco ``vazariam'' para o ambiente de
geração das instruções subsequentes, prejudicando \texttt{inferTy} e
\texttt{getVarTy}.

A função \texttt{withBlock} soluciona isso salvando e restaurando
\texttt{cgEnv} ao redor da geração de cada bloco:

\begin{minted}{haskell}
withBlock :: CgM () -> CgM ()
withBlock m = do
  saved <- gets cgEnv
  m
  modify $ \s -> s { cgEnv = saved }
\end{minted}

\subsection{Geração de \texttt{while} e \texttt{if}/\texttt{else}}

\begin{minted}{haskell}
genStmt (SWhile e body) = do
  env <- gets cgEnv
  emit ("while (" ++ genExp env e ++ ") {")
  withBlock $ indented (mapM_ genStmt body)
  emit "}"

genStmt (SIf e b1 b2) = do
  env <- gets cgEnv
  emit ("if (" ++ genExp env e ++ ") {")
  withBlock $ indented (mapM_ genStmt b1)
  emit "} else {"
  withBlock $ indented (mapM_ genStmt b2)
  emit "}"
\end{minted}

Cada ramo é gerado em um \texttt{withBlock} independente: variáveis declaradas
no \texttt{then} não aparecem no ambiente durante a geração do \texttt{else}.

\begin{Definition}[Novos comandos de TWhile]\label{def:twhile-stmt}
\begin{align*}
\mathcal{S}_\Gamma\llbracket \mathbf{while}\ e\ \{\ \vec{s}\ \} \rrbracket
  &= \texttt{while (}\mathcal{E}_\Gamma\llbracket e \rrbracket\texttt{) \{}\
     \mathcal{S}_{\Gamma'}\llbracket \vec{s} \rrbracket\
     \texttt{\}} \\
\mathcal{S}_\Gamma\llbracket \mathbf{if}\ e\ \{\ \vec{s}_1\ \}\ \mathbf{else}\ \{\ \vec{s}_2\ \} \rrbracket
  &= \texttt{if (}\mathcal{E}_\Gamma\llbracket e \rrbracket\texttt{) \{}\ \mathcal{S}_{\Gamma'}\llbracket \vec{s}_1 \rrbracket\ \texttt{\} else \{}\ \mathcal{S}_{\Gamma''}\llbracket \vec{s}_2 \rrbracket\ \texttt{\}}
\end{align*}
onde $\Gamma'$ e $\Gamma''$ são os ambientes estendidos com as declarações de
cada bloco, restritos ao escopo do respectivo bloco.
\end{Definition}

\subsection{Exemplo: fatorial iterativo}

O programa TWhile:
\begin{minted}{text}
var n : int = 5;
var acc : int = 1;
while n > 0 {
    acc := acc * n;
    n := n - 1;
}
print acc;
\end{minted}

gera:

\begin{minted}{c}
#include <stdio.h>
#include <stdlib.h>
#include <string.h>

int main(void) {
    long n = 5L;
    long acc = 1L;
    while ((n > 0L)) {
        acc = (acc * n);
        n = (n - 1L);
    }
    printf("%ld\n", (long)(acc));
    return 0;
}
\end{minted}

% ============================================================================
\section{Geração de Código para TImp}\label{sec:codegen-timp}
% ============================================================================

TImp é a linguagem mais rica da família: acrescenta \textbf{tipos registro},
\textbf{funções de usuário} com recursão mútua e as expressões \texttt{new},
\texttt{e.f} e \texttt{f(args)}.  A geração de código para TImp exige três
estratégias adicionais:

\begin{enumerate}
  \item Registros são mapeados para \texttt{struct}s C, com uma função
    auxiliar \texttt{make\_R} por tipo para alocação e inicialização.
  \item Funções de usuário são geradas com \emph{forward declarations}
    para suportar recursão mútua.
  \item O estado da mônada é estendido com dois ambientes adicionais:
    campos dos registros (\texttt{cgRecEnv}) e tipos de retorno das funções
    (\texttt{cgFnEnv}).
\end{enumerate}

\subsection{Estado estendido}

\begin{minted}{haskell}
data CgSt = CgSt
  { cgLines  :: [String]
  , cgIndent :: Int
  , cgEnv    :: Map Var Ty           -- variáveis locais em escopo
  , cgRecEnv :: Map Name [FieldDecl] -- campos de cada record
  , cgFnEnv  :: Map Name RetTy       -- tipo de retorno de cada função
  }
\end{minted}

\texttt{cgRecEnv} é necessário para gerar \texttt{ENew} na ordem dos campos
da declaração e para realizar \texttt{inferTy} em \texttt{EField}.
\texttt{cgFnEnv} é necessário para \texttt{inferTy} em chamadas de função
usadas dentro de \texttt{print}.

\subsection{Records: structs e funções construtoras}

Cada declaração \texttt{record R \{ $f_1$ : $\tau_1$; \ldots \}} gera duas
construções C:

\paragraph{1.\ Typedef struct.}

\begin{minted}{c}
typedef struct {
    long f1;
    /* demais campos */
} R;
\end{minted}

\begin{minted}{haskell}
genRecordTypedef :: RecordDecl -> CgM ()
genRecordTypedef (RecordDecl rname fields) = do
  emit "typedef struct {"
  indented $ mapM_ (\(FieldDecl f t) -> emit (cTy t ++ " " ++ f ++ ";")) fields
  emit ("} " ++ rname ++ ";")
\end{minted}

\paragraph{2.\ Função construtora \texttt{make\_R}.}

A expressão \texttt{new R \{ \ldots \}} precisa ser compilada como uma
\emph{expressão} C — sem statements auxiliares — porque pode aparecer em
posição de argumento ou de inicializador.  A solução é gerar uma função
auxiliar que aloca o struct via \texttt{malloc} e o inicializa com os valores
dos campos:

\begin{minted}{c}
R* make_R(long f1, ...) {
    R* _r = (R*)malloc(sizeof(R));
    _r->f1 = f1;
    /* demais campos */
    return _r;
}
\end{minted}

Com isso, \texttt{new R \{ $f_1$ = $e_1$, \ldots \}} compila para
\texttt{make\_R($\mathcal{E}\llbracket e_1 \rrbracket$, \ldots)}, uma chamada
de função pura sem extensões de compilador.

\begin{minted}{haskell}
genMakeHelper :: RecordDecl -> CgM ()
genMakeHelper (RecordDecl rname fields) = do
  let params = commas [cTy t ++ " " ++ f | FieldDecl f t <- fields]
  emit (rname ++ "* make_" ++ rname ++ "(" ++ params ++ ") {")
  indented $ do
    emit (rname ++ "* _r = (" ++ rname ++ "*)malloc(sizeof(" ++ rname ++ "));")
    mapM_ (\(FieldDecl f _) -> emit ("_r->" ++ f ++ " = " ++ f ++ ";")) fields
    emit "return _r;"
  emit "}"
\end{minted}

\subsection{Geração de expressões (monádica)}

Em TImp, a geração de expressões é \emph{monádica} (retorna \texttt{CgM
String}) porque consulta \texttt{cgRecEnv} e \texttt{cgFnEnv}:

\begin{minted}{haskell}
genExp :: Exp -> CgM String
genExp (EField (EVar v) f) = pure (v ++ "->" ++ f)
genExp (EField e f) = do
  ce <- genExp e
  pure (parens ce ++ "->" ++ f)
genExp (ECall fname args) = do
  cArgs <- mapM genExp args
  pure (fname ++ "(" ++ commas cArgs ++ ")")
genExp (ENew rname initFields) = do
  recEnv <- gets cgRecEnv
  case Map.lookup rname recEnv of
    Nothing   -> throwError ("Tipo record desconhecido: " ++ rname)
    Just flds -> do
      args  <- mapM (lookupField initFields rname) flds
      cArgs <- mapM genExp args
      pure ("make_" ++ rname ++ "(" ++ commas cArgs ++ ")")
\end{minted}

A tradução de \texttt{ENew} percorre os campos na \emph{ordem da declaração},
garantindo que a assinatura de \texttt{make\_R} seja sempre respeitada.

\begin{Definition}[Expressões adicionais de TImp]\label{def:timp-exp}
\begin{align*}
\mathcal{E}\llbracket x.f \rrbracket
  &= x \texttt{->} f \\
\mathcal{E}\llbracket e.f \rrbracket
  &= \texttt{(}\mathcal{E}\llbracket e\rrbracket\texttt{)->}f \\
\mathcal{E}\llbracket f(\vec{e}) \rrbracket
  &= f\texttt{(}\mathcal{E}\llbracket e_1\rrbracket\texttt{,\ \ldots,\ }\mathcal{E}\llbracket e_n\rrbracket\texttt{)} \\
\mathcal{E}\llbracket \mathbf{new}\ R\ \{\ f_1 = e_1,\ \ldots\ \} \rrbracket
  &= \texttt{make\_}R\texttt{(}\mathcal{E}\llbracket e_{\sigma(1)}\rrbracket\texttt{,\ \ldots)}
\end{align*}
onde $\sigma$ reordena os campos conforme a declaração de $R$.
\end{Definition}

\subsection{Comandos adicionais de TImp}

\begin{minted}{haskell}
genStmt (SFieldAssign v f e) = do
  ce <- genExp e
  emit (v ++ "->" ++ f ++ " = " ++ ce ++ ";")

genStmt (SReturn Nothing)  = emit "return;"
genStmt (SReturn (Just e)) = do { ce <- genExp e; emit ("return " ++ ce ++ ";") }

genStmt (SCall fname args) = do
  cArgs <- mapM genExp args
  emit (fname ++ "(" ++ commas cArgs ++ ");")
\end{minted}

\begin{Definition}[Comandos adicionais de TImp]\label{def:timp-stmt}
\begin{align*}
\mathcal{S}\llbracket v.f \mathrel{:=} e \rrbracket
  &= v\texttt{->}f\ \texttt{=}\ \mathcal{E}\llbracket e\rrbracket\texttt{;} \\
\mathcal{S}\llbracket \mathbf{return}\ e \rrbracket
  &= \texttt{return}\ \mathcal{E}\llbracket e\rrbracket\texttt{;} \qquad
\mathcal{S}\llbracket \mathbf{return} \rrbracket = \texttt{return;} \\
\mathcal{S}\llbracket f(\vec{e}) \rrbracket
  &= f\texttt{(}\mathcal{E}\llbracket \vec{e}\rrbracket\texttt{);}
\end{align*}
\end{Definition}

\subsection{Funções de usuário e forward declarations}

TImp permite funções mutuamente recursivas.  Em C, uma função só pode ser
chamada após sua declaração.  A solução é emitir \textbf{declarações avançadas}
de todas as funções antes de suas definições:

\begin{minted}{haskell}
funcSignature :: FuncDecl -> String
funcSignature (FuncDecl fname params retTy _) =
  cRetTy retTy ++ " " ++ fname ++ "(" ++ paramList ++ ")"
  where
    paramList | null params = "void"
              | otherwise   = commas [cParamTy t ++ " " ++ v | Param v t <- params]

genForwardDecl :: FuncDecl -> CgM ()
genForwardDecl fd = emit (funcSignature fd ++ ";")

genFuncDef :: FuncDecl -> CgM ()
genFuncDef fd@(FuncDecl _ params _ body) = do
  emit (funcSignature fd ++ " {")
  withBlock $ indented $ do
    mapM_ (\(Param v t) -> setVar v t) params
    mapM_ genStmt body
  emit "}"
\end{minted}

O corpo da função é gerado dentro de \texttt{withBlock}, de modo que os
parâmetros e variáveis locais não poluem o ambiente global após o retorno.

\subsection{Estrutura do código gerado para TImp}

O ponto de entrada \texttt{compileTImp} orquestra a geração na seguinte
ordem, respeitando as dependências de declaração do C:

\begin{enumerate}
  \item Cabeçalhos (\texttt{\#include}).
  \item \texttt{typedef struct \{ \ldots \} R;} para cada record.
  \item Funções \texttt{make\_R(\ldots)} para cada record.
  \item Forward declarations de todas as funções de usuário.
  \item Definições completas de todas as funções de usuário.
  \item \texttt{int main(void) \{ \ldots \}} com os comandos de topo do programa.
\end{enumerate}

\subsection{Exemplo: registros com funções}

O programa TImp:
\begin{minted}{text}
record Point { x : int; y : int; }

fn distance(p: Point): int {
    return p.x * p.x + p.y * p.y;
}

var p : Point = new Point { x = 3, y = 4 };
print distance(p);
\end{minted}

gera:

\begin{minted}{c}
#include <stdio.h>
#include <stdlib.h>
#include <string.h>

typedef struct {
    long x;
    long y;
} Point;

Point* make_Point(long x, long y) {
    Point* _r = (Point*)malloc(sizeof(Point));
    _r->x = x;
    _r->y = y;
    return _r;
}

long distance(Point * p);

long distance(Point * p) {
    return ((p->x * p->x) + (p->y * p->y));
}

int main(void) {
    Point *p = make_Point(3L, 4L);
    printf("%ld\n", (long)(distance(p)));
    return 0;
}
\end{minted}

Compilado e executado, o programa imprime \texttt{25} ($3^2 + 4^2 = 25$).

% ============================================================================
\section{Uso do Transpilador}\label{sec:transp-uso}
% ============================================================================

Os três transpiladores são acessíveis via a opção \texttt{--c}:

\begin{minted}{bash}
$ cabal run tline  -- --c -f prog.tline   # gera prog.c
$ cabal run twhile -- --c -f prog.twhile  # gera prog.c
$ cabal run timp   -- --c -f prog.timp    # gera prog.c
\end{minted}

O arquivo \texttt{.c} gerado pode então ser compilado com qualquer compilador
C compatível com C99:

\begin{minted}{bash}
$ gcc -O2 -o prog prog.c
$ ./prog
\end{minted}

% ============================================================================
\section{Comparação com a Abordagem via IRT}\label{sec:transp-vs-irt}
% ============================================================================

É instrutivo comparar as duas rotas de compilação disponíveis para TImp:

\begin{center}
\begin{tabular}{lll}
\hline
\textbf{Aspecto} & \textbf{Via IRT $\to$ WAT} & \textbf{Transpilação $\to$ C} \\
\hline
Estágios             & 4 (AST, IRT, WAT, Wasm) & 2 (AST, C) \\
Complexidade         & Alta                    & Baixa \\
Otimizações          & Manual (IRT)            & Delegadas ao \texttt{gcc} \\
Código gerado        & Texto WAT / binário Wasm & Texto C legível \\
Portabilidade        & Qualquer runtime WASI   & Qualquer plataforma com \texttt{gcc} \\
Depuração            & Requer ferramentas Wasm & Direta no .c com \texttt{gdb} \\
\hline
\end{tabular}
\end{center}

A rota via IRT é mais adequada para estudar otimizações e geração de código
de baixo nível.  A transpilação é preferível quando o objetivo é portabilidade
imediata e simplicidade de implementação — exatamente o compromisso que
TypeScript faz com JavaScript.

% ============================================================================
\section{Conclusão}\label{sec:transp-conclusao}
% ============================================================================

Este capítulo apresentou o projeto e a implementação de transpiladores de
TLine, TWhile e TImp para C.

Os pontos principais foram:

\begin{itemize}
  \item A \textbf{transpilação} é uma estratégia amplamente utilizada na
    indústria — TypeScript, CoffeeScript, Kotlin/JS, Babel e o \texttt{cfront}
    original do C++ são todos transpiladores — que delega a geração de código
    de máquina a um compilador da linguagem-alvo.

  \item O \textbf{mapeamento de tipos} exige decisões concretas: inteiros como
    \texttt{long}, booleanos como \texttt{int}, strings como buffers fixos
    \texttt{char[4096]} locais ou \texttt{char *} em parâmetros, e records como
    ponteiros para structs alocados via \texttt{malloc}.

  \item A \textbf{mônada de geração} \texttt{CgM} combina estado acumulador
    com possibilidade de falha, e é estendida progressivamente de TLine (estado
    mínimo) a TImp (com \texttt{cgRecEnv} e \texttt{cgFnEnv} adicionais).

  \item O padrão \textbf{\texttt{withBlock}} garante que o ambiente de tipos
    reflita o escopo léxico correto durante a geração, evitando que variáveis
    de blocos internos poluam o ambiente global.

  \item As \textbf{funções construtoras \texttt{make\_R}} permitem que
    \texttt{new R \{ \ldots \}} seja compilado para uma expressão C pura,
    sem extensões de compilador.

  \item As \textbf{forward declarations} de todas as funções habilitam
    recursão mútua sem exigir ordem específica de definição.
\end{itemize}
